% rel_chapter_6_sec55-56.tex
% Agent 23: Relativistic spherical perturbation equations (Ch VI §55-56)
% Conventions: (-,+,+,+) signature, c kept explicit, see BDNK_CONVENTIONS.md
% BDNK formalism: Bemfica-Disconzi-Noronha-Kovtun first-order causal theory

\section{Introduction: Spherical Geometry in the Relativistic Context}
\label{sec:rel-55}

In Chapter~VI of the classical treatment, the onset of thermal instability
was studied in spheres and spherical shells subject to a Newtonian
gravitational field.  The motivations were largely geophysical---the
convective motions in the Earth's liquid core, the pattern of mantle
convection, and related questions concerning the origin of the Earth's
magnetic field.

In the relativistic regime the spherical geometry acquires new, and arguably
more compelling, motivation.  The astrophysical objects for which general
relativity is indispensable---neutron stars, proto-neutron stars formed in
core-collapse supernovae, and the dense stellar interiors approaching
gravitational compactness---are naturally spherical (or nearly so).  In
such systems:
\begin{enumerate}
\item[(i)] The gravitational field is strong: the compactness parameter
  $\mathcal{C} = GM/(Rc^{2})$ can reach values of order $0.1$--$0.3$,
  so that Newtonian gravity is inadequate for even a qualitative
  description.
\item[(ii)] The background equilibrium is governed by the
  Tolman--Oppenheimer--Volkoff (TOV) equation rather than by classical
  hydrostatic balance.
\item[(iii)] Thermal energy can become a non-negligible fraction of the
  rest-mass energy, and the equation of state must satisfy the causality
  bound $c_{\mathrm{s}}^{2} \le c^{2}$.
\item[(iv)] Heat propagation must be treated within a causal dissipative
  framework.  We adopt the BDNK (Bemfica--Disconzi--Noronha--Kovtun)
  first-order theory~\cite{BemficaDN2018,Kovtun2019,BemficaDN2019}, in which
  causality and strong hyperbolicity are achieved through appropriate
  choice of the hydrodynamic frame and first-order constitutive relations,
  without introducing auxiliary relaxation equations.  The classical
  Fourier heat-conduction law, being parabolic, produces superluminal
  signal propagation and must be replaced.
\item[(v)] The gravitational redshift modifies the effective buoyancy
  criterion: the classical Schwarzschild--Ledoux criterion acquires a
  relativistic Tolman factor $e^{\Phi/c^{2}}$, where $\Phi$ is the
  gravitational potential appearing in the metric.
\end{enumerate}

\noindent
The programme of this and the following sections is to re-derive the
perturbation equations of Chapter~VI in a fully relativistic setting,
using the TOV equilibrium as the unperturbed background, and to
identify the corrections that arise from spacetime curvature.

%------------------------------------------------------------------
\section{The Perturbation Equations in Spherical Coordinates:
  Relativistic Treatment}
\label{sec:rel-56}

\subsection{The background spacetime and TOV equilibrium}
\label{subsec:rel-56-tov}

We consider a static, spherically symmetric spacetime.  The most general
metric of this form can be written as
\begin{equation}
  \dd s^{2}
  = -e^{2\Phi(r)/c^{2}}\,c^{2}\,\dd t^{2}
    + e^{2\Lambda(r)}\,\dd r^{2}
    + r^{2}\bigl(\dd\theta^{2} + \sin^{2}\!\theta\;\dd\varphi^{2}\bigr),
  \label{eq:rel6-metric}
\end{equation}
where the metric functions $\Phi(r)$ and $\Lambda(r)$ are determined by
Einstein's field equations.  It is conventional to introduce the
\emph{mass function}
\begin{equation}
  e^{-2\Lambda(r)}
  = 1 - \frac{2\,G\,m(r)}{r\,c^{2}} \,,
  \label{eq:rel6-Lambda}
\end{equation}
so that $m(r)$ represents the gravitational mass contained within a sphere
of areal radius~$r$.

For a perfect fluid with energy density~$\varepsilon(r)$ and
pressure~$p(r)$, Einstein's equations reduce to the
Tolman--Oppenheimer--Volkoff (TOV) system:
\begin{equation}
  \frac{\dd m}{\dd r}
  = \frac{4\pi\,r^{2}\,\varepsilon}{c^{2}} \,,
  \label{eq:rel6-mass}
\end{equation}
\begin{equation}
  \frac{\dd p}{\dd r}
  = -\frac{(\varepsilon + p)}{c^{2}}
    \,\frac{G\,m + 4\pi\,G\,r^{3}\,p/c^{2}}
          {r\bigl(r - 2\,G\,m/c^{2}\bigr)} \,,
  \label{eq:rel6-TOV}
\end{equation}
\begin{equation}
  \frac{\dd\Phi}{\dd r}
  = -\frac{c^{2}}{\varepsilon + p}\,\frac{\dd p}{\dd r} \,.
  \label{eq:rel6-Phi}
\end{equation}
Equation~\eqref{eq:rel6-TOV} is the relativistic generalization of the
classical hydrostatic equation $\dd p/\dd r = -\rho\,g(r)$, to which it
reduces in the limit $p \ll \varepsilon \approx \rho_{0}\,c^{2}$ and
$2Gm/(rc^{2}) \ll 1$.  The Newtonian gravitational acceleration
$g(r) = Gm(r)/r^{2}$ is replaced by the effective gravitational
acceleration
\begin{equation}
  g_{\mathrm{eff}}(r)
  = \frac{G\,m + 4\pi\,G\,r^{3}\,p/c^{2}}
        {r^{2}\bigl(1 - 2Gm/(rc^{2})\bigr)} \,,
  \label{eq:rel6-geff}
\end{equation}
which incorporates the contributions of pressure to the gravitational
source (a purely relativistic effect) and the spatial curvature
correction factor $\bigl(1 - 2Gm/(rc^{2})\bigr)^{-1}$.

%------------------------------------------------------------------
\subsection{Relativistic thermal background}
\label{subsec:rel-56-thermal}

In the classical treatment (eq.~6-3 of the original text), the
temperature distribution was determined by the Fourier equation
$\kappa\,\nabla^{2}T = -\epsilon$.  In the relativistic setting
we employ the BDNK first-order causal heat
equation~\cite{BemficaDN2018,Kovtun2019}.  In the BDNK framework, the
heat flux is given by the first-order constitutive relation
\begin{equation}
  Q^{\mu} = -\kappa_{\mathrm{rel}}\,T\,\Delta^{\mu\alpha}
    \partial_{\alpha}\!\left(\frac{\mu}{T}\right)
    + \text{(frame correction terms)},
  \label{eq:rel6-BDNK-heat}
\end{equation}
where the frame correction terms---specific to the BDNK general
frame---ensure that the resulting system of conservation laws is
strongly hyperbolic and hence causal, without the need for any
relaxation equation.

For a static background the time-derivative terms vanish and the
steady-state temperature profile satisfies
\begin{equation}
  \frac{1}{r^{2}}\frac{\dd}{\dd r}\!\left(
    r^{2}\,e^{-\Lambda}\,\kappa_{\mathrm{rel}}
    \,e^{-\Lambda}\,\frac{\dd(e^{\Phi/c^{2}}\,T)}{\dd r}
  \right)
  = -\epsilon_{\mathrm{rel}}\,e^{2\Lambda} \,,
  \label{eq:rel6-heat-bg}
\end{equation}
where $\kappa_{\mathrm{rel}}$ is the relativistic thermal conductivity
and $\epsilon_{\mathrm{rel}}$ is the proper volumetric heat-generation
rate.  The factor $e^{\Phi/c^{2}}\,T$ is the Tolman-redshifted
temperature, which is spatially uniform in global thermal equilibrium.
The temperature gradient driving buoyancy is therefore
\begin{equation}
  \beta_{\mathrm{rel}}(r)
  = -e^{-\Lambda}\frac{\dd T}{\dd r}
    + \frac{T}{c^{2}}\,e^{-\Lambda}\frac{\dd\Phi}{\dd r} \,,
  \label{eq:rel6-beta}
\end{equation}
where the second term is the \emph{Tolman correction}, absent in the
Newtonian theory.  In the non-relativistic limit
($\Phi/c^{2} \to 0$, $\Lambda \to 0$), $\beta_{\mathrm{rel}}$
reduces to the classical gradient
$\beta = -\dd T/\dd r = 2\bar{\beta}_{0}\,r$.

\begin{quote}
\textbf{Remark (BDNK vs.\ Israel--Stewart).}
In the Israel--Stewart (IS) second-order theory, causal heat transport is
achieved by promoting the heat flux to a dynamical variable obeying a
relaxation equation with relaxation time $\tau_{q}$.  The BDNK theory
achieves causality at first order through the choice of hydrodynamic frame;
no relaxation times appear.  For the static background considered here,
both formulations yield the \emph{same} steady-state equation
\eqref{eq:rel6-heat-bg}, since all time-derivative terms (including the IS
relaxation terms) vanish identically.  The distinction becomes relevant only
for the time-dependent perturbation equations (see
\S\,\ref{subsec:rel-56-modes}).
\end{quote}

%------------------------------------------------------------------
\subsection{The operator $\mathscr{L}^{2}$ in curved spacetime}
\label{subsec:rel-56-L2}

The angular momentum operator $\mathscr{L}^{2}$ was introduced in the
classical theory in Section~\ref{sec:56a}.  In the metric
\eqref{eq:rel6-metric} the two-sphere at fixed $r$ and $t$ has the
standard round metric $r^{2}\,\dd\Omega^{2}$, and the Laplacian on this
two-sphere is
\begin{equation}
  \mathscr{L}^{2}
  = -\frac{1}{\sin\theta}\frac{\partial}{\partial\theta}\!\left(
      \sin\theta\,\frac{\partial}{\partial\theta}
    \right)
    - \frac{1}{\sin^{2}\!\theta}\frac{\partial^{2}}{\partial\varphi^{2}} \,,
  \label{eq:rel6-L2}
\end{equation}
which is \emph{identical in form} to the classical expression
(eq.~\eqref{eq:6-25}).  This is because the angular sector of the
metric~\eqref{eq:rel6-metric} is conformally identical to flat space;
the curvature of the spacetime is encoded entirely in the
$(t,r)$-sector through $\Phi$ and $\Lambda$.  The eigenvalue equation
\begin{equation}
  \mathscr{L}^{2}\,Y_{l}^{m}(\theta,\varphi)
  = l(l+1)\,Y_{l}^{m}(\theta,\varphi)
  \label{eq:rel6-L2-eigen}
\end{equation}
therefore carries over unchanged, and the spherical harmonics
$Y_{l}^{m}(\theta,\varphi)$ remain the natural angular basis.

The full covariant Laplacian acting on a scalar function
$f(r)\,Y_{l}^{m}(\theta,\varphi)$ in the metric
\eqref{eq:rel6-metric} becomes
\begin{equation}
  \Box_{\mathrm{spatial}}\bigl[f(r)\,Y_{l}^{m}\bigr]
  = \left(\mathscr{D}_{l,\mathrm{rel}}^{2}\,f\right) Y_{l}^{m} \,,
  \label{eq:rel6-scalar-lap}
\end{equation}
where the relativistic radial operator is
\begin{equation}
  \mathscr{D}_{l,\mathrm{rel}}^{2}
  = e^{-2\Lambda}\!\left(
      \frac{\dd^{2}}{\dd r^{2}}
      + \left[\frac{2}{r} - \frac{\dd\Lambda}{\dd r}
              + \frac{1}{c^{2}}\frac{\dd\Phi}{\dd r}\right]
        \frac{\dd}{\dd r}
    \right)
    - \frac{l(l+1)}{r^{2}} \,.
  \label{eq:rel6-Dlrel}
\end{equation}
In the Newtonian limit ($\Lambda \to 0$, $\Phi/c^{2}\to 0$)
this reduces to the classical operator
$\mathscr{D}_{l}^{2} = \dd^{2}/\dd r^{2}
  + (2/r)\,\dd/\dd r - l(l+1)/r^{2}$
(cf.\ eq.~\eqref{eq:6-34}).

%------------------------------------------------------------------
\subsection{Normal mode analysis with spherical harmonics}
\label{subsec:rel-56-modes}

Following the procedure of Section~\ref{sec:56b}, we decompose
perturbations into normal modes.  In the relativistic setting we write,
for Eulerian perturbations about the TOV background,
\begin{equation}
  \delta u_{r}\,r
  = W(r)\,Y_{l}^{m}(\theta,\varphi)\,e^{pt} \,,
  \label{eq:rel6-W}
\end{equation}
\begin{equation}
  \delta T
  = \Theta(r)\,Y_{l}^{m}(\theta,\varphi)\,e^{pt} \,,
  \label{eq:rel6-Theta}
\end{equation}
\begin{equation}
  \frac{\delta p}{\varepsilon + p}
  = \Pi(r)\,Y_{l}^{m}(\theta,\varphi)\,e^{pt} \,,
  \label{eq:rel6-Pi}
\end{equation}
where $p$ (not to be confused with pressure) is the growth rate.
The key distinction from the classical analysis is that the
background quantities $\varepsilon$, $p$, and $T$ are determined by the
TOV system \eqref{eq:rel6-mass}--\eqref{eq:rel6-Phi} rather than by
Newtonian hydrostatics.

Inserting these modal forms into the linearized relativistic
conservation equations $\nabla_{\mu}(\delta T^{\mu\nu}) = 0$ with
the BDNK first-order constitutive relations, and using the angular
decomposition via $\mathscr{L}^{2}\,Y_{l}^{m} = l(l+1)\,Y_{l}^{m}$,
we obtain the coupled radial system.

\paragraph{Radial momentum (relativistic Euler):}
\begin{equation}
  p\,W
  = -r\,e^{-2\Lambda}\frac{\dd\Pi}{\dd r}
    + \frac{g_{\mathrm{eff}}}{c^{2}}\,\alpha_{\mathrm{rel}}\,r\,\Theta
    + \nu_{\mathrm{rel}}\,\mathscr{D}_{l,\mathrm{rel}}^{2}\,W
    + \mathcal{R}_{W} \,,
  \label{eq:rel6-W-eq}
\end{equation}
where $\alpha_{\mathrm{rel}}$ is the relativistic thermal expansion
coefficient,
$\nu_{\mathrm{rel}} = \shearvisc\,c^{2}/(\varepsilon+p)$ is the
relativistic kinematic viscosity, and $\mathcal{R}_{W}$ collects
curvature-coupling terms proportional to $\dd\Phi/\dd r$ and
$\dd\Lambda/\dd r$ that vanish in the Newtonian limit.

\paragraph{Thermal perturbation (BDNK first-order):}
In the BDNK framework, the heat flux obeys a first-order constitutive
relation rather than a relaxation equation.  The linearized thermal
perturbation equation takes the form
\begin{equation}
  \left(\mathscr{D}_{l,\mathrm{rel}}^{2}
        - \frac{p}{\kappa_{\mathrm{rel}}}\right)\Theta
  = -\frac{\beta_{\mathrm{rel}}}{r}\,W
    + \mathcal{R}_{\Theta} \,,
  \label{eq:rel6-Theta-eq}
\end{equation}
where $\kappa_{\mathrm{rel}}$ is the relativistic thermal diffusivity.  Causality
is guaranteed by the BDNK frame conditions on the transport
coefficients (see \S\,\ref{subsec:rel-56-summary} below), rather than
by a relaxation prefactor.  In the non-relativistic limit this reduces
to an equation of the classical form~\eqref{eq:6-35}.

\begin{quote}
\textbf{Comparison with Israel--Stewart.}
In IS theory, equation~\eqref{eq:rel6-Theta-eq} would carry an
additional prefactor $(1 + \tau_{q}\,p)$ on the left-hand side, where
$\tau_{q}$ is the heat-flux relaxation time.  This prefactor raises the
order of the characteristic polynomial and introduces an extra,
purely decaying ``relaxation mode'' that has no direct physical
counterpart.  In the BDNK theory such spurious modes are absent: the
dispersion relation contains only the physical hydrodynamic
modes~\cite{BemficaDN2018,Kovtun2019}.  Crucially, at the marginal state
$p = 0$ both formulations yield \emph{identical} equations, so all
onset-of-instability results (critical Rayleigh numbers, cell
patterns, etc.) are the same.
\end{quote}

\paragraph{Combined equation:}
In direct analogy with equation~\eqref{eq:6-23}, we can eliminate
$\Theta$ to obtain a single equation for~$W$:
\begin{equation}
  \left(\mathscr{D}_{l,\mathrm{rel}}^{2}
          - \frac{p}{\kappa_{\mathrm{rel}}}\right)\!
  \left(\mathscr{D}_{l,\mathrm{rel}}^{2}
          - \frac{p}{\nu_{\mathrm{rel}}}\right)\!
  \mathscr{D}_{l,\mathrm{rel}}^{2}\,W
  = -\frac{\psi_{\mathrm{rel}}}{\kappa_{\mathrm{rel}}\,
     \nu_{\mathrm{rel}}}\,l(l+1)\,
     \frac{\beta_{\mathrm{rel}}}{r}\,W \,,
  \label{eq:rel6-combined}
\end{equation}
where $\psi_{\mathrm{rel}}(r) = \alpha_{\mathrm{rel}}\,
g_{\mathrm{eff}}(r)/r$ is the relativistic analogue of the function
$\psi(r) = \alpha\gamma(r)$ defined in eq.~\eqref{eq:6-12}.  Unlike
the Israel--Stewart formulation, there is no additional prefactor
$(1 + \tau_{q}\,p)$; the characteristic polynomial in $p$ has the
\emph{same degree} as the classical Fourier case, with no spurious
relaxation modes.  Causality of all modes is ensured by the BDNK
conditions on the transport coefficients~\cite{BemficaDN2018,Kovtun2019,
HoultKovtun2020}.

In the non-relativistic limit
($c \to \infty$, $\Lambda \to 0$, $\Phi \to 0$)
equation~\eqref{eq:rel6-combined} reduces precisely to
equation~\eqref{eq:6-23}.

%------------------------------------------------------------------
\subsection{Boundary conditions}
\label{subsec:rel-56-bc}

Solutions of the radial system
\eqref{eq:rel6-W-eq}--\eqref{eq:rel6-Theta-eq} must satisfy boundary
conditions that are the relativistic counterparts of those in
Section~\ref{sec:56c}.

\subsubsection{Regularity at the centre ($r = 0$)}
For a complete fluid sphere (as in a neutron star), all perturbation
variables must be regular at $r = 0$.  From the structure of the
relativistic radial operator~\eqref{eq:rel6-Dlrel} and the requirement
that the metric be smooth at the origin ($\Lambda(0) = 0$,
$m(0) = 0$), we obtain the behaviour
\begin{equation}
  W(r) \sim r^{l+1}, \qquad
  \Theta(r) \sim r^{l}, \qquad
  \Pi(r) \sim r^{l},
  \qquad\text{as } r \to 0 \,.
  \label{eq:rel6-bc-centre}
\end{equation}
This is the direct relativistic analogue of the classical condition
(eq.~\eqref{eq:6-48}) that $W$ and $G$ behave like $r^{l}$ near the
origin; the shift by one power in $W$ arises from the definition
$\delta u_{r}\,r = W\,Y_{l}^{m}\,e^{pt}$.

\subsubsection{Matching conditions at the stellar surface ($r = R$)}
At the stellar surface, defined by $p(R) = 0$, the interior solution
must match smoothly to the exterior Schwarzschild vacuum.  The
conditions are:
\begin{enumerate}
\item[(i)] The Lagrangian pressure perturbation vanishes:
  \begin{equation}
    \Delta p\big|_{r=R}
    = \delta p + \xi^{r}\frac{\dd p}{\dd r}\bigg|_{r=R} = 0 \,,
    \label{eq:rel6-bc-surface-p}
  \end{equation}
  where $\xi^{r}$ is the radial displacement.

\item[(ii)] The induced metric on the surface and its first derivative
  are continuous (junction conditions).  In terms of the perturbation
  variables this requires continuity of $W$ and $\dd W/\dd r$ across
  $r = R$.

\item[(iii)] The temperature perturbation satisfies either
  $\Theta(R) = 0$ (isothermal surface) or a radiative boundary
  condition matching to an optically thin envelope.

\item[(iv)] In the BDNK framework, the heat flux $Q^{\mu}$ is
  determined algebraically by the first-order constitutive relations;
  no additional dynamical boundary condition on $Q^{\mu}$ is required
  beyond the continuity of the fundamental variables
  $(\varepsilon, p, u^{\mu}, T)$.  This is in contrast to the
  Israel--Stewart formulation, where the dynamical heat-flux
  variable $q^{\mu}$ would require an independent boundary condition.
\end{enumerate}

\noindent
For a spherical shell confined between $r = R_{1}$ (inner boundary) and
$r = R_{0}$ (outer boundary), conditions analogous to the classical ones
(eq.~\eqref{eq:6-38}) apply at both boundaries:
\begin{equation}
  W = 0, \quad \Theta = 0
  \qquad\text{at } r = R_{1} \text{ and } r = R_{0} \,.
  \label{eq:rel6-bc-shell}
\end{equation}
For a rigid boundary the additional condition is
$\dd W/\dd r = 0$ (cf.\ eq.~\eqref{eq:6-43}); for a free boundary,
$\dd^{2}W/\dd r^{2} = 0$ (cf.\ eq.~\eqref{eq:6-47}), both now
evaluated in the coordinate $r$ of the metric~\eqref{eq:rel6-metric}.

%------------------------------------------------------------------
\subsection{Velocity field: toroidal--poloidal decomposition in
  curved spacetime}
\label{subsec:rel-56-TP}

The classical decomposition of an arbitrary solenoidal velocity field into
toroidal (\textbf{T}) and poloidal (\textbf{S}) components
(Section~\ref{sec:56d}) extends naturally to the relativistic setting.
In a static, spherically symmetric spacetime
the spatial metric on a $t = \text{const}$ slice takes the form
\begin{equation}
  \dd\ell^{2}
  = e^{2\Lambda}\,\dd r^{2}
    + r^{2}\,\dd\Omega^{2} \,,
  \label{eq:rel6-spatial}
\end{equation}
and the spatial part of the 4-velocity perturbation $\delta u^{i}$ can
be decomposed on each two-sphere as
\begin{equation}
  \delta u^{i}
  = \delta u^{r}\,\hat{e}_{r}^{i}
    + \frac{1}{r}\left(
        \nabla_{\!\perp}^{i}\,\mathscr{S}
        + \epsilon^{ij}_{\perp}\,\nabla_{\!\perp\,j}\,\mathscr{T}
      \right) e^{pt} \,,
  \label{eq:rel6-TP-decomp}
\end{equation}
where $\nabla_{\!\perp}^{i}$ is the covariant derivative on the
unit two-sphere, $\epsilon^{ij}_{\perp}$ is the Levi-Civita tensor on
the two-sphere, and $\mathscr{S}(r,\theta,\varphi)$,
$\mathscr{T}(r,\theta,\varphi)$ are the poloidal and toroidal scalar
potentials, respectively.

The poloidal and toroidal fields retain the fundamental properties
catalogued in the classical case (eqs.~\eqref{eq:6-49}--\eqref{eq:6-57}):
\begin{enumerate}
\item[(i)] Both fields are solenoidal with respect to the spatial
  covariant derivative $D_{i}$ of the metric~\eqref{eq:rel6-spatial}.

\item[(ii)] The orthogonality relations
  $\int \mathbf{S}\cdot\mathbf{S}'\,\dd\Omega = 0$ (for different
  harmonics), $\int \mathbf{T}\cdot\mathbf{T}'\,\dd\Omega = 0$, and
  $\int \mathbf{S}\cdot\mathbf{T}\,\dd\Omega = 0$ continue to hold,
  because these integrals are performed on the round two-sphere whose
  geometry is unaffected by the $(t,r)$-curvature.

\item[(iii)] The curl relations between poloidal and toroidal fields
  acquire metric factors: $\operatorname{curl}\mathbf{S}$ yields a
  toroidal field whose defining scalar involves the factor
  $e^{-\Lambda}$, and similarly for
  $\operatorname{curl}\mathbf{T}$.

\item[(iv)] The defining scalars $\mathscr{S}$ and $\mathscr{T}$ are
  related to the radial functions $W$ and $Z$ (the radial vorticity
  component) by
  \begin{equation}
    W = -\frac{l(l+1)}{r}\,\mathscr{S}\,, \qquad
    Z = -\frac{l(l+1)}{r}\,\mathscr{T}\,,
    \label{eq:rel6-WZ-ST}
  \end{equation}
  in direct correspondence with eq.~\eqref{eq:6-58}.
\end{enumerate}

\noindent
An important consequence for the thermal instability problem is that,
just as in the classical theory, the \emph{toroidal velocity field
is not driven by thermal buoyancy}: the buoyancy force is purely radial
and excites only the poloidal component.  The toroidal modes decouple
from the temperature perturbation and, in the absence of rotation or
magnetic fields, decay through viscous dissipation.  This structure
persists in the relativistic theory because the buoyancy force remains
radial in a spherically symmetric background.

%------------------------------------------------------------------
\subsection{BDNK causality conditions for spherical perturbations}
\label{subsec:rel-56-causality}

In the BDNK framework, causality of the perturbation equations is
guaranteed by requiring that the transport coefficients satisfy
coupled nonlinear inequalities that ensure strong
hyperbolicity~\cite{BemficaDN2018,Kovtun2019,BemficaDNK2023}.
Specifically:
\begin{enumerate}
\item[(i)] Positive shear viscosity: $\eta > 0$.
\item[(ii)] Positive bulk-plus-shear combination:
  $\zeta + \tfrac{2}{3}\eta > 0$.
\item[(iii)] The BDNK frame coefficients (the energy correction
  $\varepsilon_{1}$, etc.) satisfy a set of coupled inequalities that
  bound all characteristic speeds by~$c$.
\end{enumerate}
These conditions replace the Israel--Stewart requirement that the
relaxation times $\tau_{\pi}$, $\tau_{q}$, $\tau_{\Pi}$ be positive.
In the BDNK theory, the characteristic speeds of the linearized system
are:
\begin{itemize}
\item \textbf{Sound modes:} Modified by viscosity, with speed $\leq c$.
\item \textbf{Shear modes:}
  $\omega = -i\,\eta\,k^{2}/[(\varepsilon + p)/c^{2}] + O(k^{3})$
  (diffusive at low $k$, causal at high $k$).
\item \textbf{Heat mode:} Similarly bounded by~$c$.
\end{itemize}
Unlike the IS theory, there are no additional ``relaxation modes''---the
dispersion relation contains only the standard hydrodynamic
modes~\cite{Kovtun2019,HoultKovtun2020}.

%------------------------------------------------------------------
\subsection{Summary of relativistic corrections}
\label{subsec:rel-56-summary}

The principal modifications introduced by the relativistic treatment
of the spherical perturbation equations may be summarized as follows:
\begin{enumerate}
\item The Newtonian hydrostatic equilibrium is replaced by the TOV
  system \eqref{eq:rel6-mass}--\eqref{eq:rel6-Phi}.  All background
  quantities ($\varepsilon$, $p$, $T$, $g_{\mathrm{eff}}$) reflect the
  relativistic equation of state and the contribution of pressure to
  gravity.

\item The radial differential operator $\mathscr{D}_{l}^{2}$
  (eq.~\eqref{eq:6-34}) is replaced by
  $\mathscr{D}_{l,\mathrm{rel}}^{2}$ (eq.~\eqref{eq:rel6-Dlrel}),
  which incorporates the metric functions $\Phi$ and $\Lambda$.

\item The angular operator $\mathscr{L}^{2}$ and the spherical-harmonic
  decomposition are unchanged, since the angular sector of the metric
  is the standard round two-sphere.

\item The Fourier heat equation is replaced by the BDNK first-order
  causal constitutive relation~\cite{BemficaDN2018,Kovtun2019}.
  Unlike the Israel--Stewart second-order approach, the BDNK formulation
  does not introduce relaxation times or auxiliary dynamical variables;
  causality is ensured by the choice of hydrodynamic frame.  The
  perturbation equations have the same differential order as the
  classical (Fourier) case, with no spurious relaxation modes.

\item The temperature gradient driving buoyancy acquires a
  Tolman-redshift correction (eq.~\eqref{eq:rel6-beta}).

\item Boundary conditions at the stellar surface include junction
  conditions to the exterior Schwarzschild metric.  In the BDNK
  formulation, no additional boundary condition on the heat flux is
  needed (the heat flux is determined algebraically from the
  first-order constitutive relation), in contrast to the IS theory
  where continuity of the dynamical $q^{\mu}$ would be required.

\item The effective gravitational acceleration
  $g_{\mathrm{eff}}$ (eq.~\eqref{eq:rel6-geff}) includes pressure
  contributions to the gravitational source and a spatial-curvature
  factor.
\end{enumerate}

\noindent
All relativistic corrections vanish in the limit
$c \to \infty$ (equivalently $GM/(Rc^{2}) \to 0$),
and the classical equations of Chapter~VI are recovered exactly.
At the marginal state ($p = 0$), the BDNK and IS formulations yield
identical equations, so all critical Rayleigh numbers and onset
criteria are independent of the choice of causal framework.

%------------------------------------------------------------------
\subsection{Astrophysical application: neutron star convection setup}
\label{subsec:rel-56-astro}

The formalism developed above finds its primary astrophysical
application in neutron star interiors, where convective instability
may be driven by thermal or compositional gradients during cooling
or accretion episodes.  The essential prerequisite for any stability
analysis is the construction of the unperturbed TOV equilibrium
profile---the relativistic generalisation of the classical hydrostatic
background.

For a typical neutron star with $M \approx 1.4\,M_{\odot}$ and
$R \approx 10$~km, the compactness is $\mathcal{C} \approx 0.21$,
placing the system firmly in the regime where the TOV corrections
are quantitatively significant.  The Schwarzschild interior solution
(uniform energy density) provides a useful analytic benchmark:
figure~\ref{fig:tov-profiles} displays the radial profiles of
energy density, pressure, and the metric potential $\Phi(r)/c^{2}$
for several compactness values spanning the range from mildly
relativistic ($\mathcal{C} = 0.05$) to strongly compact
($\mathcal{C} = 0.30$).

\begin{figure}[ht]
\centering
\includegraphics[width=\textwidth]{fig_tov_profiles.pdf}
\caption{TOV equilibrium profiles for the Schwarzschild interior
  (uniform energy density) solution at four compactness values
  $\mathcal{C} = GM/(Rc^{2})$.
  \emph{Left:} Energy density (constant) and enthalpy density
  $(\varepsilon+p)/\varepsilon_{c}$.
  \emph{Centre:} Pressure profile $p(r)/\varepsilon_{c}$.
  \emph{Right:} Metric potential $\Phi(r)/c^{2}$.
  The steepening of $p(r)$ and deepening of the potential well
  with increasing compactness control the relativistic corrections
  to convective stability
  (cf.\ \S\,\ref{subsec:rel-56-summary}).}
\label{fig:tov-profiles}
\end{figure}

Several features are noteworthy for the stability analysis:
\begin{enumerate}
\item[(i)] The central pressure rises steeply with compactness
  (diverging at the Buchdahl limit $\mathcal{C} = 4/9$), so
  the enthalpy--inertia factor $\mathcal{I} = 1 + p/\varepsilon$
  is substantially greater than unity in the interior.

\item[(ii)] The metric potential $\Phi(r)$ is more negative in
  the interior for larger $\mathcal{C}$, producing a stronger
  Tolman redshift correction to the temperature
  gradient~\eqref{eq:rel6-beta}.

\item[(iii)] The effective gravitational acceleration
  $g_{\mathrm{eff}}$~\eqref{eq:rel6-geff} receives both a
  pressure-source contribution and a curvature factor
  $(1-2Gm/(rc^{2}))^{-1}$, both of which amplify buoyancy
  driving in compact stars relative to the Newtonian estimate.
\end{enumerate}

These profiles form the background on which the perturbation
equations~\eqref{eq:rel6-W-eq}--\eqref{eq:rel6-Theta-eq} are
solved.  Realistic nuclear equations of state (APR, SLy, BSk;
see~\cite{Lattimer2001,HaenselPotekhinYakovlev2007}) produce non-uniform
$\varepsilon(r)$ profiles, but the qualitative structure of the
relativistic corrections is captured by the uniform-density model
shown here.

\subsection{TOV profiles for three nuclear equations of state}
\label{subsec:rel-56-tov-3eos}

While the Schwarzschild interior solution (uniform density)
provides analytic insight, quantitative modelling of neutron star
convection requires TOV solutions for realistic nuclear equations
of state.  We solve the TOV system
\eqref{eq:rel6-mass}--\eqref{eq:rel6-Phi} numerically for three
representative EOSs using single-polytrope approximations
($p = K\,\rho^{\Gamma}$, $\varepsilon = \rho\,c^{2} +
p/(\Gamma - 1)$) calibrated to the nuclear-density
behaviour~\cite{AkmalPR1998,DouchinHaensel2001,PotekhinEtAl2013}:

\begin{center}
\begin{tabular}{lcccc}
\hline
EOS & $\Gamma$ & $R$ (km) & $M/M_{\odot}$ & $\mathscr{C}$ \\
\hline
APR  & 2.58 & $\sim 10$   & $\sim 1.4$ & $\sim 0.21$ \\
SLy  & 2.35 & $\sim 8$    & $\sim 1.2$ & $\sim 0.23$ \\
BSk21 & 2.60 & $\sim 9$   & $\sim 1.6$ & $\sim 0.29$ \\
\hline
\end{tabular}
\end{center}

Figure~\ref{fig:tov-3eos-profiles} displays the resulting radial
profiles of energy density $\varepsilon(r)/\varepsilon_{c}$,
pressure $p(r)/p_{c}$, and the pressure-to-energy ratio
$\xi(r) = p(r)/\varepsilon(r)$ inside the star.

\begin{figure}[ht]
\centering
\includegraphics[width=\textwidth]{fig_tov_3eos_profiles.pdf}
\caption{TOV profiles for three nuclear EOSs (APR, SLy, BSk21),
  solved numerically with polytropic approximations.
  \emph{Left:} Normalised energy density
  $\varepsilon(r)/\varepsilon_{c}$.  Unlike the uniform-density
  (Schwarzschild interior) model, realistic EOSs produce
  monotonically decreasing $\varepsilon(r)$.
  \emph{Centre:} Normalised pressure $p(r)/p_{c}$.  Stiffer EOSs
  (APR) have more slowly declining pressure profiles.
  \emph{Right:} Pressure-to-energy ratio
  $\xi(r) = p(r)/\varepsilon(r)$, which controls the
  enthalpy--inertia factor $\mathcal{I} = 1 + \xi$ in the
  perturbation equations.  Dotted gray lines show the
  Schwarzschild interior comparison.  The central value
  $\xi_{c} \sim 0.1$--$0.3$ quantifies the relativistic
  enhancement of buoyancy driving.}
\label{fig:tov-3eos-profiles}
\end{figure}

Key observations for the convective stability analysis:
\begin{enumerate}
\item[(i)] The energy density $\varepsilon(r)$ drops to zero at
  the surface (unlike the uniform-density model), so the
  perturbation integrals weight the inner region more heavily.

\item[(ii)] The central $\xi_{c}$ ranges from $\sim 0.06$ (low
  compactness) to $\sim 0.32$ (high compactness), directly
  controlling the magnitude of the enthalpy--inertia correction
  $\bar{\Xi}$ in the critical Rayleigh number formula
  (equation~\eqref{eq:rel-59-CrelvsC} of \S\,\ref{sec:rel-59}).

\item[(iii)] The $\xi(r)$ profile is approximately linear near
  the centre and drops more steeply near the surface, consistent
  with the ansatz $\Xi(r) \approx \Xi_{c}(1 - r^{2}/R^{2})$
  used in the post-Newtonian expansion.

\item[(iv)] The stiffer EOS (APR, larger $\Gamma$) produces a
  broader pressure distribution, while softer EOSs concentrate
  the pressure more toward the centre, affecting the radial
  structure of the convection cells at onset.
\end{enumerate}

%======================================================================
\section{Radiation Hydrodynamic Dispersion Relations and Convective Stability
  in Spherical Geometry}
\label{sec:rel-rad-hydro-dispersion}
%======================================================================

In radiation-dominated stellar interiors---particularly in massive stars
where the ratio of radiation pressure to gas pressure
$\mathcal{R} = P_{\mathrm{rad}}/P_{\mathrm{gas}} \sim T^3/n$ can reach
order unity---the thermal photon gas participates actively in the fluid
dynamics.  The convective stability analysis of this chapter must account
for the fact that thermal transport is mediated by radiation with finite
mean free path $\tau$, leading to a rich dispersion relation structure
that interpolates between optically thick (diffusive) and optically thin
(free-streaming) regimes.

\citet{GavassinoRadDispersion2025} has computed the linearized dispersion
relations for shear, heat, and sound waves in relativistic
``matter$+$radiation'' fluids with grey absorption opacities, solving the
coupled radiation-hydrodynamic equations perturbatively in the ratio
$\mathcal{R} = T_R^{00}/T_M^{00}$.  The results are exact for all
wavenumbers $k \in \mathbb{R}$.

\subsubsection{Dispersion relations}
The three types of hydrodynamic modes have the following dispersion relations:
\begin{align}
  \text{Shear:}\quad
  \omega(k) &= -i\,\frac{15\,D_s}{2\tau^2}
  \left[\frac{2}{3} + \frac{1}{(k\tau)^2}
  - \left(1 + \frac{1}{(k\tau)^2}\right)\frac{\arctan(k\tau)}{k\tau}\right],
  \label{eq:rad-shear} \\[6pt]
  \text{Heat:}\quad
  \omega(k) &= -i\,\frac{3\,D_h}{\tau^2}
  \left[1 - \frac{\arctan(k\tau)}{k\tau}\right],
  \label{eq:rad-heat} \\[6pt]
  \text{Sound:}\quad
  \omega(k) &= c_s\,k - i\,\frac{15\,D_s}{2\tau^2}
  \left[\frac{1}{3} - \frac{1 - ic_s k\tau}{(k\tau)^2}
  + \frac{(1 - ic_s k\tau)^2}{(k\tau)^3}\,
    \arctan\!\left(\frac{k\tau}{1 - ic_s k\tau}\right)\right],
  \label{eq:rad-sound}
\end{align}
where $D_s = 4\pi^2 T^4\tau/[225(\varepsilon + P)]$ and
$D_h = 4\pi^2 T^3\tau/(45\,n\,c_p)$ are the radiation-induced
diffusivities for shear and heat respectively, and $c_s$ is the
adiabatic sound speed of the matter component.

\subsubsection{Key properties}
\begin{enumerate}
  \item \textbf{Optically thick limit} ($k\tau \ll 1$): the dispersion
    relations reduce to the standard diffusive forms
    $\omega \approx -i\,D_s\,k^2$ (shear), $\omega \approx -i\,D_h\,k^2$
    (heat), recovering the Rosseland-mean opacity predictions.
  \item \textbf{Optically thin limit} ($k\tau \gg 1$): the damping rate
    saturates at a finite value $\sim D_s/\tau^2$ or $\sim D_h/\tau^2$,
    reflecting the fact that radiation on scales smaller than a mean free
    path freely streams away without further increasing the damping.
  \item \textbf{Covariant stability}: all modes satisfy
    $\mathrm{Im}\,\omega \le 0$ for real $k$, and the covariant
    stability condition $\mathrm{Im}\,\omega \le |\mathrm{Im}\,k|$
    holds for all $k \in \mathbb{C}$, ensuring the dispersion relations
    are compatible with relativistic causality.
\end{enumerate}

\subsubsection{Implications for spherical convective stability}
For the spherical perturbation problem studied in this chapter, the
radiation-hydrodynamic dispersion relations modify the thermal
diffusivity entering the critical Rayleigh number.  In the optically
thick limit, the effective thermal diffusivity is
$\kappa_{\text{eff}} = D_h = 4\pi^2 T^3 \tau/(45\,n\,c_p)$, which
scales as $T^3\tau$---the familiar Rosseland result.  The radiation
pressure also contributes to the effective adiabatic index through
$\Gamma_{\text{eff}} = [\Gamma_{\text{gas}} + 4\mathcal{R}(4\mathcal{R}+1)^{-1}]/(1+\mathcal{R})$,
modifying the Schwarzschild criterion for convective stability.

The transition from optically thick to optically thin regimes occurs
at a wavelength $\lambda_* \sim 2\pi\tau$, which for radiation in
NS envelopes ($\tau \sim 0.01$--$1\;\mathrm{cm}$) falls well within
the convective scale.  This means that the highest spherical harmonic
modes $\ell$ probed by convection may be damped more efficiently than
predicted by the diffusion approximation, effectively truncating the
convective spectrum at high $\ell$.

\begin{thebibliography}{99}
\bibitem{Lattimer2001}
J.\,M.~Lattimer and M.~Prakash,
``Neutron star structure and the equation of state,''
\emph{Astrophys.\ J.} \textbf{550}, 426--442 (2001).

\bibitem{HaenselPotekhinYakovlev2007}
P.~Haensel, A.\,Y.~Potekhin, and D.\,G.~Yakovlev,
\emph{Neutron Stars~1: Equation of State and Structure},
Springer (2007).

\bibitem{OppenheimerVolkoff1939}
J.\,R.~Oppenheimer and G.\,M.~Volkoff,
``On massive neutron cores,''
\emph{Phys.\ Rev.} \textbf{55}, 374--381 (1939).

\bibitem{Tolman1939}
R.\,C.~Tolman,
``Static solutions of Einstein's field equations for spheres of fluid,''
\emph{Phys.\ Rev.} \textbf{55}, 364--373 (1939).

\bibitem{GavassinoRadDispersion2025}
L.~Gavassino,
``Dispersion relations of relativistic radiation hydrodynamics,''
arXiv:2412.00275 (2024).

\bibitem{GavassinoCausalRadTransfer2025}
L.~Gavassino,
``Causality constraints on radiative transfer,''
arXiv:2502.08740 (2025).

\bibitem{Spiegel1957}
E.\,A.~Spiegel,
``The smoothing of temperature fluctuations by radiative transfer,''
\emph{Astrophys.\ J.} \textbf{126}, 202 (1957).
\end{thebibliography}
